\documentclass[11pt]{article}
% Shared preamble for the AmberMeta SOP and its companion note.
%
% Deliberately confined to what ships with a distro TeX Live (no tlmgr on the
% workstation these are built on): every package below is in texlive-latex-recommended
% or texlive-fonts-recommended.
\usepackage[T1]{fontenc}
\usepackage[utf8]{inputenc}
\usepackage{newtxtext,newtxmath}
\usepackage[scaled=0.85]{inconsolata}
\usepackage{microtype}
\usepackage[a4paper,margin=25mm,headheight=14pt]{geometry}
\usepackage{graphicx}
\usepackage{booktabs}
\usepackage{tabularx}
\usepackage{enumitem}
\usepackage{titlesec}
\usepackage{fancyhdr}
\usepackage{xcolor}
\usepackage{tcolorbox}
\usepackage{caption}
\usepackage{url}
\usepackage[hidelinks]{hyperref}

\definecolor{amber}{HTML}{B45309}
\definecolor{ink}{HTML}{1F2937}
\definecolor{muted}{HTML}{6B7280}
\definecolor{rule}{HTML}{D1D5DB}
\definecolor{panel}{HTML}{F9FAFB}

\color{ink}

% -- headings ------------------------------------------------------------------
\titleformat{\section}{\large\bfseries\color{amber}}{\thesection}{0.6em}{}
\titleformat{\subsection}{\normalsize\bfseries\color{ink}}{\thesubsection}{0.6em}{}
\titlespacing*{\section}{0pt}{1.4em}{0.5em}
\titlespacing*{\subsection}{0pt}{1.0em}{0.35em}

% -- lists ---------------------------------------------------------------------
\setlist[itemize]{leftmargin=1.2em,itemsep=0.25em,topsep=0.35em}
\setlist[enumerate]{leftmargin=1.5em,itemsep=0.45em,topsep=0.4em}

% -- captions ------------------------------------------------------------------
\captionsetup{font=small,labelfont={bf,color=amber},skip=6pt}

% -- code ----------------------------------------------------------------------
% \code is for inline literals. Underscores are escaped at the call site; that is
% verbose but predictable, and it keeps paths copy-pasteable out of the PDF.
\newcommand{\code}[1]{\texttt{\small #1}}

% Literal paths that may break across lines inside a narrow table column. Breaks after
% `/` like \path does, and takes underscores unescaped.
\DeclareUrlCommand\bpath{\urlstyle{tt}\small}

\newtcolorbox{shell}{
  colback=panel, colframe=rule, boxrule=0.4pt, arc=2pt,
  left=6pt, right=6pt, top=5pt, bottom=5pt, fontupper=\ttfamily\small,
}

\newtcolorbox{note}[1]{
  colback=panel, colframe=amber, boxrule=0pt, leftrule=2pt, arc=0pt,
  left=8pt, right=8pt, top=6pt, bottom=6pt,
  title={\bfseries\color{amber}\small #1}, coltitle=amber,
  fonttitle=\bfseries\small, attach title to upper={\par\smallskip},
}

% -- page furniture ------------------------------------------------------------
\pagestyle{fancy}
\fancyhf{}
\renewcommand{\headrulewidth}{0.4pt}
\renewcommand{\footrulewidth}{0pt}
\fancyfoot[C]{\small\color{muted}\thepage}


\fancyhead[L]{\small\color{muted}Replica layouts in the deposited corpus}
\fancyhead[R]{\small\color{muted}AmberMeta v1.1.0}

\begin{document}

\begin{center}
  {\LARGE\bfseries\color{amber} How Replicas Are Organised, and What\\[0.2em]
   AmberMeta Can Detect}\\[0.7em]
  {\large Findings from the deposited simulation corpus}\\[1.1em]
  {\small Michele Bonus \quad\textbullet\quad AmberMeta v1.1.0 \quad\textbullet\quad \today}
\end{center}
\vspace{0.6em}
\hrule
\vspace{1.3em}

\section{Scope}

This note accompanies the AmberMeta SOP and answers the standing question of how users
organise replicas and whether AmberMeta can recognise them. It reports what was found in
the deposited corpus at \code{/store7/gentile/data/simulations}, what AmberMeta does with
each layout today, and what would have to change to cover the cases it declines.

No code was changed. Everything below was measured against AmberMeta v1.1.0 as released.

\section{The corpus}

Of 50 system directories, 19 hold simulation files. 30 are empty placeholders. One,
\code{sys005}, is \emph{unreadable}: its permissions are \code{drw-rw-r-{}-}, so the
execute bit is missing and no user (including its owner) can traverse it. That is a
one-line fix at the source (\code{chmod~u+rx}) and worth doing before any further
deposits.

\subsection{What was tested}

Five systems were selected to cover every replica convention present. All were scanned
read-only; nothing was written into the corpus.

\begin{center}\small
\begin{tabularx}{\textwidth}{@{}lrrrr>{\raggedright\arraybackslash}X@{}}
\toprule
\textbf{System} & \textbf{Files} & \textbf{Logs} & \textbf{Steps} & \textbf{Scan} & \textbf{Layout}\\
\midrule
\code{sys023} & 1044 & 205  & 213  & 6.4\,s  & Chunked, no replicas\\
\code{sys011} & 1280 & 260  & 268  & 1.9\,s  & \code{run1..run4} beside the stages\\
\code{sys004} & 5100 & \textbf{0} & 107  & 83.5\,s & \code{run1..run6} inside each stage\\
\code{sys021} & 5348 & 1091 & 1105 & 27.8\,s & \code{01..05} inside each stage\\
\code{sys024} & 167  & 30   & 50   & 1.3\,s  & \code{L12Y/prod/01..05}\\
\bottomrule
\end{tabularx}
\end{center}

\medskip
All five were parsed without error. Two results stand out independently of replicas:

\begin{itemize}
  \item \textbf{\code{sys004} contains no output logs at all.} 5100 files (2496
        restarts, 2496 trajectories, 104 inputs) and not one \code{.mdout}. Less is lost
        than that sounds. The trajectories are NetCDF despite their \code{.mdcrd} names,
        and AmberMeta types a file by its content, so each run still reports its atom
        count, frame count, frame spacing, achieved simulation time, box type and volume
        statistics: one production chunk here reads 5200 to 6000 ps from the trajectory
        alone. What the missing logs cost is the temperature, the completion status, the
        wall clock, and the authoritative record of where each run began. The last of
        those has a consequence of its own, in \S\ref{sec:quirk}.
  \item \textbf{\code{sys011} has 260 logs but only 65 inputs.} The runs are recoverable
        from the logs, but the requested settings for most of them are not on record.
\end{itemize}

\section{The four conventions, and what AmberMeta does}

\begin{center}\small
\begin{tabularx}{\textwidth}{@{}l>{\raggedright\arraybackslash}Xcc@{}}
\toprule
\textbf{Convention} & \textbf{Shape} & \textbf{Auto} & \textbf{Manual}\\
\midrule
Replicas within stages   & \code{equilibration/run1..run6}, \code{production/run1..run8} & yes & n/a\\
Numbered within stages   & \code{equil/01..05}, \code{prod/01..05}                        & yes & n/a\\
Replicas beside stages   & \code{run1..run4} next to \code{equilibration/equi1..equi4}    & no  & yes\\
Id repeated in filenames & \code{L12Y/prod/01/GB1\_204\_01\_000}                          & no  & yes\\
\bottomrule
\end{tabularx}
\end{center}

\medskip
Two of the four are recognised with no intervention, and produce correctly tagged
members (\code{run1..run8} for \code{sys004}; \code{01..05} for \code{sys021}). The other
two are \emph{declined}, not mis-grouped; in both, the runs are still discovered
and still correctly chained. Only the grouping into members is withheld.

\subsection{Why the two declines happen}

Neither is a defect. AmberMeta's layout inference tags nothing when a tree admits more
than one reading, on the grounds that a wrong member claim is worse than none.

\textbf{\code{sys011}} presents two families whose tag sets are disjoint:
\code{equilibration/equi1..equi4}, which varies at the second path segment, and
\code{run1..run4}, which varies at the first. Disjoint families are refused because they
are equally consistent with two separate experiments in one directory. Restricting the
scan to the replica directories alone resolves it immediately: \code{discover
-{}-pattern 'run[0-9]'} tags \code{run1..run4} correctly.

\textbf{\code{sys024}} repeats the replica identifier inside the filenames:
\code{prod/01/GB1\_204\_01\_000}, \code{prod/02/GB1\_204\_02\_000}. Grouping keys on the
set of run \emph{bases} a directory holds, so each replica directory ends up with a base
set of its own, no two directories are ever seen as running the same thing, and no cohort
with more than one member is ever formed. This one cannot be fixed by narrowing the
scan; the identifier has to be picked manually.

\subsection{The manual route works for both}

\textbf{Define replicas\ldots} lets the user choose which path segment names the member.
On \code{sys011} choosing the first segment offers
\code{equilibration}, \code{leap}, \code{run1}, \code{run2}, \code{run3}, \code{run4}; on
\code{sys024} choosing the third offers the preparation step names alongside
\code{01..05}. In both cases the replicas are present and correct among the options, and
accepting only those produces a manifest identical to one from a conforming layout.

The cost is judgement: the picker offers every directory at the chosen level, including
preparation directories that are not replicas. A user who accepts all of them gets a
manifest asserting that \code{leap} is a replica.

\subsection{A parser quirk this corpus exposed}
\label{sec:quirk}

\code{sys004}'s production runs hold a trajectory and a restart and nothing else. Files
are grouped by stem, so a run's own output restart lands in the slot continuity reads its
\emph{start} time from. v1.1.0 guards against exactly that, but the guard asks whether the
group also holds an \code{mdin} or an \code{mdout}, on the reasoning that a group with
neither is not a run at all but a bare topology and coordinate pair. Here both are absent
and the runs are real, so the guard never fires: each chunk's own restart is read as its
start time, and every junction reports a gap of one full chunk (999.9999998 ps against a
1000 ps chunk).

The equilibration directories carry \code{mdin} files, so the guard fires there and
continuity reports an honest ``cannot verify'' instead. The condition is narrow, a deposit
carrying neither inputs nor logs, but that is the shape a partial upload takes, and inside
it the phantom gap the guard exists to prevent returns in full.

\section{Recommendations}

\begin{enumerate}
  \item \textbf{Adopt the layout in \S3 of the SOP for new work.} Replicas as sibling
        directories directly inside their stage, identically named across stages, with
        the identifier in the directory name only. Both auto-detected conventions in the
        corpus already follow this; it costs nothing to standardise on it.

  \item \textbf{Treat missing output logs as a deposit blocker.} \code{sys004} is the
        case to cite. The trajectories carry the timing, so the deposit is not worthless,
        but no run in it can be shown to have finished, no temperature is on record, and
        the continuity of every production chain is reported wrongly rather than
        cautiously (\S\ref{sec:quirk}). A one-line check at deposit time, every run
        directory holding a \code{.mdout}, would have caught it.

  \item \textbf{Raise the self-restart guard's condition as an issue.} Requiring an
        \code{mdin} or an \code{mdout} in the group is what makes it miss
        \code{sys004}. Keying on whether the group's coordinate file is also the run's
        declared output would cover the log-free case, and is worth a design pass rather
        than a quick patch.

  \item \textbf{Fix \code{sys005}'s permissions} and re-run the survey over all 19
        populated systems once the SOP is in use.

  \item \textbf{Consider two contained improvements to the picker}, in order of value:
        \begin{itemize}
          \item Mark which offered members \emph{look like} replicas (directories whose
                run-base sets match each other) and which do not, so a user is steered
                away from accepting \code{leap} as a member. This is presentational and
                does not touch the inference rule.
          \item Consider whether a directory whose base sets differ only by an embedded
                copy of its own directory name (the \code{sys024} shape) can be
                recognised safely. This is a real change to the membership predicate and
                would need its own design; the safe-failure property is worth more than
                the convenience, so it should not be rushed.
        \end{itemize}

  \item \textbf{Do not weaken the disjoint-families rule.} The \code{sys011} refusal is
        the rule working as intended. The procedural workaround is cheap and the failure
        mode it prevents, silently merging two experiments into one campaign, is
        expensive.
\end{enumerate}

\end{document}
